\documentclass[11pt,a4paper]{article}

\usepackage[utf8]{inputenc}
\usepackage[T1]{fontenc}
\usepackage{geometry}
\usepackage{amsmath}
\usepackage{longtable}
\usepackage{array}
\usepackage{url}

\geometry{margin=2.5cm}

\newcommand{\toprule}{\hline}
\newcommand{\midrule}{\hline}
\newcommand{\bottomrule}{\hline}

\title{Technical Report: Fine-Tuning a Route Replanning Model}
\author{AdaptiveRoute POC}
\date{August 25, 2026}

\begin{document}

\maketitle

\begin{abstract}
This report documents the design, training, evaluation, and final selection of the specialized route replanning model used in the AdaptiveRoute proof of concept. The approach combines a prescriptive optimization layer, implemented with Pyomo and HiGHS, with supervised fine-tuning of a neural route-candidate policy. The fine-tuned model is not treated as a mathematical optimizer and does not replace the solver. It is used as a fast structured candidate generator, protected by deterministic validation, repair logic, and a Pyomo + HiGHS fallback. The selected final version is the LoRA adapter at \path{outputs/models/adaptiveroute-qwen2_5-7b-lora-error20k-v5}, based on \path{Qwen/Qwen2.5-7B-Instruct}.
\end{abstract}

\section{Executive Summary}

The fine-tuning track produced a specialized model for generating JSON route replanning candidates under operational disruptions. The final selected model, referred to as v5, reached a 94.4\% feasible-candidate rate on a fixed 1,000-example benchmark sampled from \path{data/training_curriculum_100k/sft_test.jsonl}. It improved over the earlier v3 and v4 checkpoints and outperformed the later v6 continuation, which slightly regressed.

The key architectural decision is to keep optimization and learning separated:

\begin{itemize}
  \item Pyomo + HiGHS is the prescriptive source of truth.
  \item The fine-tuned LLM is a route-candidate policy.
  \item Deterministic validation decides whether a candidate can be accepted.
  \item Repair and solver fallback remain mandatory.
\end{itemize}

This makes the system suitable for an agentic AI architecture where a reasoning-capable orchestrator can delegate route candidate generation to the fine-tuned specialist model.

\section{Project Objective}

The goal of the fine-tuning module was to train a specialist model that receives:

\begin{itemize}
  \item a structured routing scenario;
  \item an existing route plan;
  \item an operational event, such as a blocked arc or unavailable customer;
  \item vehicle capacities, customer demands, and route constraints;
\end{itemize}

and returns a replanned route candidate in a strict JSON format.

The proof of concept focuses on synthetic CVRP instances, primarily with 8 customers and 2 vehicles. The target behavior is operational feasibility, not proof of optimality. A feasible generated plan is useful only if it passes deterministic validation.

\section{Theoretical Foundation}

\subsection{Vehicle Routing and CVRP}

The Vehicle Routing Problem was introduced by Dantzig and Ramser in the classical truck dispatching formulation \cite{dantzig1959}. In the Capacitated Vehicle Routing Problem, a fleet of vehicles with limited capacity must serve a set of customers from a depot while minimizing total travel cost or distance.

CVRP is a combinatorial optimization problem. For realistic instances, exact optimization can become expensive, and practical systems often combine mathematical programming, heuristics, and learned policies. In this project, the learned model is not positioned as an exact solver. It is trained to imitate high-quality prescriptive decisions generated by a mathematical optimization engine.

\subsection{Prescriptive Optimization}

The prescriptive layer answers the question: given a scenario and a disruption, what route plan should be executed? Pyomo was selected because it allows explicit algebraic modeling in Python and keeps the optimization model inspectable, testable, and version-controlled \cite{pyomo}. HiGHS was selected as the open-source solver because it provides high-performance optimization routines for LP, MIP, and QP problems and has a Python interface suitable for local experimentation \cite{highs}.

This combination supports the intended technical message: the route optimization model is explicit and auditable, rather than hidden inside a black-box model.

\subsection{Transformer-Based Language Models}

The base model is a causal Transformer model \cite{vaswani2017}. The attention mechanism is useful for this task because the model must condition route decisions on relationships among customers, vehicles, capacity constraints, blocked arcs, and the previous plan.

An instruction-tuned model was preferred because the task is contract-driven: the model must follow a structured prompt and produce valid JSON.

\subsection{LoRA and QLoRA}

Full fine-tuning of a 7B-parameter model was not practical for the local hardware target. Instead, the project used parameter-efficient fine-tuning with LoRA. LoRA freezes the original pretrained weights and injects trainable low-rank matrices into selected layers, drastically reducing the number of trainable parameters \cite{lora}.

The training setup also used QLoRA-style 4-bit loading. QLoRA enables fine-tuning large language models by keeping the base model quantized while training LoRA adapters, using techniques such as NF4 quantization and double quantization \cite{qlora}.

Therefore, the full base model was not trained. The base \path{Qwen/Qwen2.5-7B-Instruct} model remained frozen; only LoRA adapter weights were updated.

\section{Selected Base Model}

The selected base model was \path{Qwen/Qwen2.5-7B-Instruct}. The selection criteria were practical:

\begin{itemize}
  \item open-weight model suitable for local experimentation;
  \item approximately 7.61B parameters;
  \item 28 Transformer layers;
  \item modern architecture features, including RoPE, SwiGLU, RMSNorm, and Grouped Query Attention;
  \item strong instruction-following behavior for structured output;
  \item compatibility with Hugging Face Transformers, TRL, PEFT, and bitsandbytes.
\end{itemize}

The model is intended to serve as a specialized routing policy. It should not be the same model responsible for user-facing reasoning. In the planned agentic AI architecture, a stronger reasoning model can handle conversation, orchestration, and explanation, while the fine-tuned model serves route-specialist agents or tools.

\section{Prescriptive Optimization Component}

The prescriptive engine was implemented with Pyomo + HiGHS. It has three roles:

\begin{enumerate}
  \item solve synthetic CVRP instances to generate supervised labels;
  \item validate whether LLM-generated candidates satisfy deterministic constraints;
  \item provide fallback plans when generated candidates are invalid or missing.
\end{enumerate}

A simplified CVRP objective is:

\[
\min \sum_{k \in K} \sum_{i \in N} \sum_{j \in N, j \neq i} d_{ij} x_{ijk}
\]

subject to customer coverage:

\[
\sum_{k \in K} \sum_{i \in N, i \neq j} x_{ijk} = 1
\quad \forall j \in C
\]

capacity:

\[
\sum_{j \in C} q_j y_{jk} \leq Q_k
\quad \forall k \in K
\]

flow conservation:

\[
\sum_{j \in N, j \neq i} x_{ijk}
=
\sum_{j \in N, j \neq i} x_{jik}
\quad \forall i \in N, \forall k \in K
\]

and blocked-arc exclusion:

\[
x_{ijk} = 0
\quad \forall (i,j) \in B, \forall k \in K
\]

where $N$ is the node set, $C$ is the customer set, $K$ is the vehicle set, $d_{ij}$ is travel distance, $q_j$ is customer demand, $Q_k$ is vehicle capacity, $B$ is the blocked-arc set, $x_{ijk}$ indicates whether vehicle $k$ traverses arc $(i,j)$, and $y_{jk}$ indicates whether vehicle $k$ serves customer $j$.

The implementation also includes depot start/end rules, route continuity, subtour elimination, and event-specific constraints.

\section{Dataset Generation}

\subsection{Generation Methodology}

The datasets were generated synthetically to provide scale, repeatability, and reliable labels. The generation pipeline was:

\begin{enumerate}
  \item generate a CVRP instance with depot, customers, demands, vehicles, capacities, and distance matrix;
  \item solve the base route plan with Pyomo + HiGHS;
  \item apply an operational event, such as a blocked arc or unavailable customer;
  \item re-solve the disrupted scenario with the solver;
  \item validate the replanned route;
  \item serialize the scenario, event, and target output into compact SFT JSONL;
  \item convert compact SFT rows into chat-message JSONL for instruction fine-tuning.
\end{enumerate}

This approach makes the LLM learn from prescriptive optimization outputs while preserving a deterministic oracle for auditing and runtime validation.

\subsection{Generated Datasets}

\begin{longtable}{>{\raggedright\arraybackslash}p{5.4cm}rrrr>{\raggedright\arraybackslash}p{3.5cm}}
\toprule
Dataset & Total & Train & Validation & Test & Purpose \\
\midrule
\endhead
\path{training_compact_30k} & 30000 & 24000 & 3000 & 3000 & Initial compact curriculum. \\
\path{training_capacity_tight_50k_parallel} & 50000 & 40000 & 5000 & 5000 & Emphasis on capacity pressure. \\
\path{training_blocked_arc_20k_parallel} & 20000 & 16000 & 2000 & 2000 & Emphasis on blocked arcs. \\
\path{training_curriculum_100k} & 100000 & 80000 & 10000 & 10000 & Consolidated curriculum dataset. \\
\path{training_error_driven_20k} & 20000 & 16000 & 2000 & 2000 & Continuation data focused on observed errors. \\
\path{training_error_driven_10k_v6} & 10000 & 8000 & 1000 & 1000 & Additional continuation after v5. \\
\bottomrule
\end{longtable}

The corresponding chat-format datasets were stored under \path{data/chat_training_*}. The main split strategy was 80\% training, 10\% validation, and 10\% test. Local audits found zero invalid rows in the principal datasets used for training.

\section{Fine-Tuning Strategy}

\subsection{Training Algorithm}

The selected training algorithm was supervised fine-tuning with LoRA/QLoRA. Training used:

\begin{itemize}
  \item Hugging Face Transformers for model loading;
  \item TRL \path{SFTTrainer} for supervised fine-tuning;
  \item PEFT for LoRA adapter management;
  \item bitsandbytes for 4-bit model loading.
\end{itemize}

The optimization target is causal language-model loss over the assistant response tokens. In practice, the model learns to map a structured route replanning prompt to the expected JSON route candidate.

\subsection{Final LoRA Configuration}

\begin{table}[h]
\centering
\begin{tabular}{ll}
\toprule
Parameter & Value \\
\midrule
Base model & \path{Qwen/Qwen2.5-7B-Instruct} \\
Method & LoRA / QLoRA-style loading \\
PEFT type & \path{LORA} \\
Task type & \path{CAUSAL_LM} \\
Rank $r$ & 16 \\
\path{lora_alpha} & 32 \\
\path{lora_dropout} & 0.05 \\
Bias & \path{none} \\
Target modules & \path{q_proj,k_proj,v_proj,o_proj,gate_proj,up_proj,down_proj} \\
Quantization & 4-bit NF4 with double quantization \\
Compute dtype & bfloat16 \\
\bottomrule
\end{tabular}
\end{table}

This configuration trains adapter matrices in attention and MLP projection modules. The original dense weights of the base model remain frozen.

\section{Training Methodology and Parameters}

Training was conducted in stages. Each continuation was evaluated against a fixed benchmark before deciding whether to keep training or freeze the current model.

\begin{longtable}{>{\raggedright\arraybackslash}p{1.1cm}rrrrrr>{\raggedright\arraybackslash}p{3.2cm}}
\toprule
Version & Train rows & Eval rows & Steps & Batch & LR & Time & Notes \\
\midrule
\endhead
v3 & 24000 & 500 & 3000 & 1 $\times$ 8 accum. & $5\cdot10^{-5}$ & 5h53m & First stable full training on the 30K compact dataset. \\
v4 & 56000 & 1000 & 2300 & 8 & $1.7\cdot10^{-5}$ & 3h44m & Hard-data continuation using SDPA and checkpoint resume. \\
v5 & 16000 & 1000 & 1500 & 8 & $1\cdot10^{-5}$ & 2h22m & Error-driven continuation. Selected final version. \\
v6 & 8000 & 500 & 800 & 8 & $7\cdot10^{-6}$ & 1h18m & Additional continuation; slightly regressed. \\
\bottomrule
\end{longtable}

The stabilized training setup used:

\begin{itemize}
  \item \path{--max-seq-length 2048} for v4/v5/v6 continuations;
  \item \path{--bf16};
  \item \path{--bnb-compute-dtype bfloat16};
  \item \path{--optim adamw_torch};
  \item \path{--lr-scheduler-type cosine};
  \item \path{--max-grad-norm 0.3};
  \item \path{--no-eval-during-train};
  \item \path{dataset_num_proc=1};
  \item \path{dataloader_num_workers=0}.
\end{itemize}

The representative v5 training command was:

\begin{verbatim}
uv run python scripts/train_lora.py \
  --model-id Qwen/Qwen2.5-7B-Instruct \
  --adapter-path outputs/models/adaptiveroute-qwen2_5-7b-lora-hard70k-v4-b8-sdpa-resume700 \
  --dataset-dir data/chat_training_error_driven_20k \
  --output-dir outputs/models/adaptiveroute-qwen2_5-7b-lora-error20k-v5 \
  --max-seq-length 2048 \
  --epochs 1 \
  --max-steps 1500 \
  --learning-rate 1e-5 \
  --warmup-steps 50 \
  --max-grad-norm 0.3 \
  --lr-scheduler-type cosine \
  --per-device-batch-size 8 \
  --gradient-accumulation-steps 1 \
  --train-limit 16000 \
  --eval-limit 1000 \
  --bf16 \
  --bnb-compute-dtype bfloat16 \
  --optim adamw_torch \
  --save-steps 100 \
  --logging-steps 10 \
  --no-eval-during-train
\end{verbatim}

\section{Training Feedback}

\subsection{Training Metrics}

The logs show that loss quickly converged to the 0.26--0.28 range after the initial learning stage. Mean token accuracy stabilized around 0.89. These metrics indicate that the model learned the response format and common routing patterns, but they are not sufficient to prove operational correctness. The primary evaluation criterion remained route feasibility.

\begin{table}[h]
\centering
\begin{tabular}{lrrrr}
\toprule
Version & Train loss & Final token accuracy & Tokens processed & Runtime \\
\midrule
v3 & 0.2796 & 0.8922 & $1.614\cdot10^7$ & 21190s \\
v4 & 0.2659 & 0.8915 & $1.241\cdot10^7$ & 13470s \\
v5 & 0.2654 & 0.8911 & $8.081\cdot10^6$ & 8510s \\
v6 & 0.2653 & 0.8917 & $4.307\cdot10^6$ & 4700s \\
\bottomrule
\end{tabular}
\end{table}

\subsection{Operational Evaluation}

All candidate models were evaluated on the same 1,000-example sample from \path{data/training_curriculum_100k/sft_test.jsonl}. The main metric was feasible-candidate rate. Exact match was tracked but treated as secondary because CVRP can have multiple valid route plans for the same scenario.

\begin{table}[h]
\centering
\begin{tabular}{lrrrrr}
\toprule
Model & Valid JSON & Feasible & Exact match & Capacity violations & Blocked-arc violations \\
\midrule
v3 & 100.0\% & 91.7\% & 4.6\% & 72 & 11 \\
v4 & 100.0\% & 92.9\% & 5.8\% & 61 & 10 \\
v5 & 100.0\% & 94.4\% & 5.6\% & 46 & 9 \\
v6 & 100.0\% & 94.0\% & 5.4\% & 50 & 10 \\
\bottomrule
\end{tabular}
\end{table}

The main result is that v5 improved feasibility and reduced capacity violations relative to v4. The later v6 continuation did not improve the fixed benchmark and was not selected.

\section{Limitations Found}

\subsection{Model Limitations}

\begin{itemize}
  \item The model does not guarantee mathematical feasibility or optimality.
  \item Residual violations remain. On the 1,000-example benchmark, v5 still produced 46 capacity violations and 9 blocked-arc violations.
  \item Exact match is low, although this is partly expected for CVRP because multiple feasible solutions can exist.
  \item Training data is synthetic; no real traffic, fleet telemetry, driver constraints, service times, or geospatial road-network data were used in this stage.
  \item The primary training scope was small CVRP instances, approximately 8 customers.
\end{itemize}

\subsection{Infrastructure Limitations}

\begin{itemize}
  \item Long-running generation and training jobs exposed native failures, including \path{exit status 139}, double-free errors, and a \path{torch.nn.Module.train} runtime issue.
  \item FlashAttention was not successfully installed because the local environment used Python 3.12, Torch with CUDA 13.0, and an nvcc/CUDA stack without a compatible prebuilt wheel.
  \item The stable workaround was to use SDPA, reduce preprocessing parallelism, disable evaluation during training, and use \path{adamw_torch}.
\end{itemize}

\section{Refinement Strategy}

The next refinements should target operational feasibility, not marginal loss reduction. The recommended strategy is:

\begin{enumerate}
  \item integrate v5 into the agentic routing module behind a deterministic validator;
  \item implement repair for simple capacity, missing-customer, and blocked-arc violations;
  \item use Pyomo + HiGHS fallback when validation or repair fails;
  \item generate new error-driven datasets from actual validator failures;
  \item train new continuations only if they beat v5 on the fixed benchmark;
  \item test best-of-N route candidates with deterministic selection;
  \item expand synthetic data to more customers, more vehicles, time windows, multiple depots, asymmetric costs, and richer disruption types;
  \item include public VRP/CVRP benchmarks when moving beyond the current POC scope;
  \item revisit FlashAttention or Unsloth only after stabilizing a compatible CUDA stack.
\end{enumerate}

\section{Final Version and Runtime Contract}

The selected model is:

\begin{verbatim}
outputs/models/adaptiveroute-qwen2_5-7b-lora-error20k-v5
\end{verbatim}

The stable operational alias is:

\begin{verbatim}
outputs/models/adaptiveroute-routing-policy-current
\end{verbatim}

The recommended runtime contract is:

\begin{verbatim}
scenario + base_plan + event
-> LoRA v5 candidate generation
-> JSON parsing
-> deterministic route validation
-> repair if needed
-> Pyomo + HiGHS fallback if still invalid
-> final response
\end{verbatim}

The final accepted response should include candidate source, validation result, rejected violations if any, route distance, vehicle loads, and final route JSON.

\section{Training and Test Environment}

\begin{table}[h]
\centering
\begin{tabular}{ll}
\toprule
Item & Observed value \\
\midrule
Operating system & Ubuntu 24.04.4 LTS \\
Kernel & Linux 6.8.0-138-generic \\
CPU & Intel Core i9-14900K, 32 logical CPUs \\
GPU & NVIDIA GeForce RTX 3090 \\
VRAM & 24576 MiB \\
NVIDIA driver & 580.173.02 \\
Compute capability & 8.6 \\
Python & 3.12.3 via \path{uv} \\
Torch & 2.13.0, CUDA 13.0 \\
Transformers & 5.15.1 \\
TRL & 1.10.0 \\
PEFT & 0.20.0 \\
bitsandbytes & 0.50.1 \\
datasets & 5.0.1 \\
Pyomo & 6.10.1 \\
HiGHS/highspy & 1.15.1 \\
\bottomrule
\end{tabular}
\end{table}

\section{Time Spent}

\begin{longtable}{>{\raggedright\arraybackslash}p{5cm}>{\raggedright\arraybackslash}p{4.2cm}>{\raggedright\arraybackslash}p{5cm}}
\toprule
Activity & Observed time & Notes \\
\midrule
\endhead
Initial smoke training & 3m24s & 20 steps; validated the basic pipeline. \\
v3 30K training & 5h53m & First stable full training run. \\
v4 hard 70K continuation & 3h44m & Continued from checkpoint with SDPA. \\
v5 error-driven 20K continuation & 2h22m & Selected final model. \\
v6 error-driven 10K continuation & 1h18m & Did not outperform v5. \\
1,000-example prediction generation & 20--25min per model & Observed during v5/v6 evaluation runs. \\
Synthetic dataset generation & Long chunked/parallel jobs & Monitored with \path{screen}, \path{systemd-inhibit}, and \path{tee} logs. \\
\bottomrule
\end{longtable}

\section{Conclusion}

The selected approach is technically defensible for a five-day proof of concept. The prescriptive Pyomo + HiGHS layer provides deterministic labels, validation, and fallback. The LoRA/QLoRA fine-tuning process specializes an instruction-following 7B model into a route-candidate generator. The final model decision was based on operational feasibility, not only training loss.

The v5 adapter is the current best stopping point. It achieved the highest feasible-candidate rate among evaluated models and reduced the most important violation categories. The next engineering step should be integration into the LangGraph-based agentic AI module, with a validator, repair layer, and solver fallback. Further training should be driven by runtime validation failures rather than by generating more synthetic data without a measured error target.

\begin{thebibliography}{9}

\bibitem{dantzig1959}
G. B. Dantzig and J. H. Ramser.
\newblock The Truck Dispatching Problem.
\newblock \textit{Management Science}, 6(1):80--91, 1959.
\newblock \url{https://doi.org/10.1287/mnsc.6.1.80}

\bibitem{vaswani2017}
A. Vaswani et al.
\newblock Attention Is All You Need.
\newblock arXiv:1706.03762, 2017.
\newblock \url{https://arxiv.org/abs/1706.03762}

\bibitem{lora}
E. J. Hu et al.
\newblock LoRA: Low-Rank Adaptation of Large Language Models.
\newblock arXiv:2106.09685, 2021.
\newblock \url{https://arxiv.org/abs/2106.09685}

\bibitem{qlora}
T. Dettmers, A. Pagnoni, A. Holtzman and L. Zettlemoyer.
\newblock QLoRA: Efficient Finetuning of Quantized LLMs.
\newblock arXiv:2305.14314, 2023.
\newblock \url{https://arxiv.org/abs/2305.14314}

\bibitem{qwen25}
Qwen Team.
\newblock Qwen2.5 Technical Report and Qwen2.5-7B-Instruct model card.
\newblock \url{https://arxiv.org/abs/2412.15115} and \url{https://huggingface.co/Qwen/Qwen2.5-7B-Instruct}

\bibitem{pyomo}
M. L. Bynum et al.
\newblock Pyomo -- Optimization Modeling in Python.
\newblock Third Edition, Springer, 2021.
\newblock \url{https://www.pyomo.org/citing-pyomo}

\bibitem{highs}
HiGHS.
\newblock High-performance open-source software for linear optimization.
\newblock \url{https://highs.dev/}

\end{thebibliography}

\end{document}
